\textbf{COLUNI 2026 -- Questão 33}\par\medskip Leia o trecho a seguir, extraído da obra \textit{Os Sertões} (1902), de Euclides da Cunha:

“Então, a travessia das veredas sertanejas é mais exaustiva que a de uma estepe nua. (...) Ao passo que a caatinga o afoga; abrevia-lhe o olhar; agride-o e estonteia-o; enlaça-o na trama espinescente e não o atrai; repulsa-o com as folhas urticantes, com o espinho, com os gravetos estalados em lanças (...): árvores sem folhas, de galhos estorcidos e secos (...), lembrando um bracejar imenso, de tortura, da flora agonizante...”\par
(Disponível em: \url{https://fundar.org.br/wp-content/uploads/2021/06/os-sertoes.pdf}. Adaptado. Acesso em: 30 jul. 2025.)

Com base no trecho acima e nos conhecimentos sobre o meio físico e humano do Nordeste brasileiro, analise as afirmativas abaixo:
\begin{enumerate}
\item[I.] O texto revela uma percepção estética e simbólica do domínio morfoclimático da caatinga, frequentemente associado à aridez e à hostilidade ambiental, aspectos que impactam diretamente as dinâmicas socioeconômicas da região semiárida.
\item[II.] A descrição da vegetação como “bracejar de tortura” traduz uma crítica à degradação ambiental promovida pela agricultura extensiva e pelo desmatamento recentes na caatinga.
\item[III.] A vegetação da caatinga apresenta adaptações ao clima semiárido, como folhas reduzidas, raízes profundas, espinhos e estruturas que minimizam a perda de água.
\item[IV.] Apesar da aparente “desolação” descrita no trecho, a caatinga é um dos domínios morfoclimáticos mais biodiversos do Brasil, mas sofre com a escassez de políticas públicas voltadas à sua conservação.
\end{enumerate}
Está CORRETO o que se afirma apenas em:\begin{enumerate}
\item I e III.
\item II e IV.
\item I e II.
\item I, III e IV.\end{enumerate}